\documentclass[11pt,a4paper]{article}

\usepackage[utf8]{inputenc}
\usepackage[T1]{fontenc}
\usepackage[margin=2.4cm]{geometry}
\usepackage{amsmath}
\usepackage{amssymb}
\usepackage{siunitx}
\usepackage{booktabs}
\usepackage{graphicx}
\usepackage{float}
\usepackage{caption}
\usepackage{xcolor}
\usepackage[hidelinks]{hyperref}

\graphicspath{{results/}}
\sisetup{group-separator={\,},detect-weight=true}
\DeclareSIUnit{\veh}{veh}

\captionsetup{font=small,labelfont=bf}

\title{\textbf{Urban Traffic Controller}\\[2pt]
\large Large-scale optimisation of traffic-signal timings\\
with a parallel genetic algorithm on Cyfronet Ares}
\author{Grzegorz Mróz \and Katarzyna Kucharczyk}
\date{June 2026}

\begin{document}
\maketitle

\tableofcontents
\bigskip

\begin{abstract}
\noindent
We present an end-to-end system that designs road networks, optimises their
traffic-signal timings, and scales the optimisation across a supercomputer.
A microscopic traffic simulator (SUMO) scores a candidate signal plan by the
\emph{total vehicle delay} it produces; a genetic algorithm (GA) searches the
space of green-phase durations to minimise that delay. Because every fitness
evaluation is an independent SUMO run, the population is evaluated in parallel
with MPI, turning a sequential search into a high-performance computing (HPC)
workload. On a single $8\times8$ city the optimised plan cuts total vehicle
delay by \SI{27.3}{\percent} versus fixed-time signals. We empirically justify
the choice of a GA by comparing it, on an equal evaluation budget, against
random search, hill climbing and simulated annealing --- the GA reaches the
deepest improvement. On Cyfronet Ares the parallel evaluator delivers a
\textbf{$26.96\times$ speedup on 47 workers} and sustains up to
\SI{45.8}{evaluations\per\second}, while a problem-size study shows the cost per
evaluation growing gracefully from a 9-intersection grid to a 1024-intersection
city. An interactive dashboard ties the pieces together: build a network,
optimise it, and watch the before/after traffic.
\end{abstract}
\bigskip

% ====================================================================
\section{Problem statement}
% ====================================================================
Traffic lights at an intersection cycle through a fixed sequence of phases.
The duration of each \emph{green} phase decides how much of the cycle a given
movement receives. Poorly chosen durations force vehicles to wait at red while
a conflicting approach sits empty, wasting time and fuel across the whole
network. The classical engineering answer --- fixed-time plans from rules of
thumb --- ignores the actual demand pattern of a specific city.

We treat signal timing as a black-box optimisation problem. Let a city contain
intersections $i=1,\dots,N$, each with $g_i$ controllable green phases. A
\emph{plan} is the vector of all green durations,
\[
  \mathbf{x} = (x_1,\dots,x_n), \qquad n = \sum_{i=1}^{N} g_i,
  \qquad x_j \in [\,g_{\min}, g_{\max}\,] \subset \mathbb{Z}.
\]
Yellow and all-red phases stay fixed, so every plan we generate is a valid,
safe signal program. The quality of a plan is the \emph{total vehicle delay}
it produces over a simulated horizon $T$:
\begin{equation}
  D(\mathbf{x}) \;=\; \sum_{v \in V} \int_{0}^{T}
     \mathbf{1}\!\left[\,\text{vehicle } v \text{ is halting at time } t\,\right]\,dt,
  \label{eq:delay}
\end{equation}
i.e.\ the cumulative vehicle-seconds spent stopped. We seek the plan
$\mathbf{x}^\star = \arg\min_{\mathbf{x}} D(\mathbf{x})$. The objective has no
closed form: $D(\mathbf{x})$ is only available by running a full microscopic
traffic simulation, which makes the problem a natural fit for population-based,
derivative-free search --- and an expensive one, motivating parallelism.

% ====================================================================
\section{Tools and methods}
% ====================================================================
\begin{description}
  \item[SUMO] (Simulation of Urban MObility) --- an open-source microscopic
  traffic simulator. Each vehicle is modelled individually; the simulator
  reports per-step halting counts from which Eq.~\eqref{eq:delay} is computed.
  Networks are generated with \texttt{netgenerate} and demand with
  \texttt{randomTrips.py}. On the cluster SUMO is driven in-process through the
  \texttt{libsumo} bindings, avoiding socket overhead.
  \item[DEAP] --- a Python evolutionary-computation framework providing the GA
  operators (selection, crossover, mutation) we configure for the integer
  genome.
  \item[mpi4py] --- MPI bindings for Python. We use
  \texttt{MPIPoolExecutor} from \texttt{mpi4py.futures} to evaluate the
  population in a master--worker pattern.
  \item[Dash / Plotly] --- the interactive dashboard front-end: build a
  network, launch the optimiser, and visualise the before/after traffic.
  \item[Cyfronet Ares] --- the HPC platform (partition \texttt{plgrid}). SUMO
  v1.20.0 is built from source with \texttt{libsumo}; OpenMPI 4.1.6 launches the
  ranks via PMIx.
\end{description}

% ====================================================================
\section{System architecture}
% ====================================================================
The system is a single pipeline that can run serially on a laptop or in
parallel on the cluster without code changes:
\[
\boxed{\text{Dashboard}} \;\rightarrow\;
\boxed{\text{Genetic algorithm}} \;\rightarrow\;
\boxed{\text{SUMO simulation}}
\]
The dashboard builds or uploads a network; the GA proposes signal plans; SUMO
scores each plan. Crucially, the GA dispatches every evaluation through an
abstract \texttt{map\_fn}. Locally this is Python's built-in \texttt{map}
(serial); on the cluster it is \texttt{MPIPoolExecutor.map} (parallel). The
\emph{same} optimiser code therefore drives both the single-machine demo and the
HPC scaling study --- only the map function changes.

% ====================================================================
\section{Problem formulation}
% ====================================================================
\paragraph{Genome.} A candidate plan is encoded as an integer vector: one gene
per controllable green phase, holding its duration in seconds and bounded by
$[g_{\min},g_{\max}]=[\SI{5}{\second},\SI{60}{\second}]$. For an $8\times8$ grid
this yields $n=128$ genes; for a $32\times32$ city, $n=2048$. Because only green
durations are encoded and the protective yellow/red phases are untouched, the
decoder always emits a legal SUMO signal program.

\paragraph{Fitness.} Decoding a genome writes the durations into the network's
signal programs; one headless SUMO run then returns the total vehicle delay of
Eq.~\eqref{eq:delay}. Lower is better. A single \emph{baseline} run with the
default fixed-time plan provides the reference against which every result is
reported as an improvement percentage.

% ====================================================================
\section{The optimisation algorithm}
% ====================================================================
\subsection{Why a genetic algorithm}
The objective is non-differentiable, noisy and multimodal: small changes in one
intersection's timing interact non-linearly with traffic arriving from its
neighbours. There are no usable gradients. A genetic algorithm is well suited
because it (i) needs only the scalar fitness, (ii) maintains a \emph{population}
of plans, exploring many regions of the search space at once, and (iii) exposes
exactly the parallelism we need --- the whole population is evaluated
independently each generation. Section~\ref{sec:compare} backs this choice with
an empirical comparison against simpler search strategies.

\subsection{Operators}
\begin{description}
  \item[Selection] tournament selection (size 3): pick the fitter of randomly
  drawn plans, biasing reproduction towards good solutions while keeping
  diversity.
  \item[Crossover] uniform crossover: each child gene is taken from either
  parent with equal probability, recombining good per-intersection timings.
  \item[Mutation] per-gene uniform reset within $[g_{\min},g_{\max}]$, which
  injects fresh durations and prevents premature convergence.
  \item[Elitism] a Hall of Fame of size~1 carries the best plan unchanged into
  the next generation, so the optimum never regresses --- the result can never
  be worse than the baseline.
\end{description}

\subsection{Pseudocode}
\noindent\hrulefill
\begin{verbatim}
initialise population P of random integer plans
evaluate D(x) for every x in P        # via map_fn -> SUMO (parallel)
record best-so-far in Hall of Fame
for generation = 1 .. G:
    offspring <- tournament_select(P)
    offspring <- uniform_crossover(offspring, p_cx)
    offspring <- uniform_reset_mutation(offspring, p_mut)
    evaluate D(x) for every new x       # via map_fn -> SUMO (parallel)
    P <- offspring + elite(Hall of Fame)
    update Hall of Fame
return best plan in Hall of Fame
\end{verbatim}
\noindent\hrulefill

\medskip
The only line that touches the cluster is \texttt{evaluate} --- it submits the
generation's plans through \texttt{map\_fn}. Everything else (selection,
recombination, bookkeeping) is cheap and runs on the master.

% ====================================================================
\section{Algorithm comparison: why the GA earns its place}
\label{sec:compare}
% ====================================================================
A reviewer may reasonably ask whether the GA's machinery is worth it, or whether
a simpler search would do as well. To answer this empirically we ran four
algorithms on the \emph{same} objective (the $8\times8$ city, baseline delay
\SI{39473}{\veh\second}) and gave each the \emph{same evaluation budget} --- the
GA's own evaluation count ($\approx1990$ SUMO runs). All metaheuristics are
seeded with the baseline plan (mirroring the GA's elitism), so none can report a
result worse than fixed-time; the question is purely \emph{how far below the
baseline each can reach on an equal budget}.

\begin{description}
  \item[Random search] samples random plans and keeps the best.
  \item[Hill climbing] parallel random-restart local search: each chain accepts
  a one-gene neighbour only if it improves.
  \item[Simulated annealing] like hill climbing, but accepts worsening moves
  with probability $e^{-\Delta/T}$ under a geometric cooling schedule.
  \item[Genetic algorithm] the production optimiser described above.
\end{description}

\begin{table}[H]
\centering
\caption{Final result of each algorithm on an equal evaluation budget
($8\times8$ city, baseline $D=\SI{39473}{\veh\second}$). Lower delay is better.}
\label{tab:compare}
\begin{tabular}{lS[table-format=5.0]S[table-format=2.1]S[table-format=4.0]}
\toprule
Algorithm & {Best delay [veh$\cdot$s]} & {Improvement [\%]} & {Evaluations} \\
\midrule
\textbf{Genetic algorithm} & \textbf{28689} & \textbf{27.3} & 1991 \\
Hill climbing              & 31693          & 19.7          & 1978 \\
Random search              & 36010          & 8.8           & 1991 \\
Simulated annealing        & 36405          & 7.8           & 1978 \\
\midrule
Baseline (fixed-time)      & 39473          & 0.0           & 1 \\
\bottomrule
\end{tabular}
\end{table}

\begin{figure}[H]
\centering
\includegraphics[width=\textwidth]{algorithms.png}
\caption{Algorithm comparison on equal budget. \emph{Left:} convergence ---
best total delay against the number of SUMO evaluations; the GA descends below
every competitor. \emph{Right:} final improvement over the fixed-time baseline.
The GA reaches \SI{-27.3}{\percent}, hill climbing \SI{-19.7}{\percent}, while
random search and simulated annealing stall near \SI{-8}{\percent}.}
\label{fig:compare}
\end{figure}

\paragraph{Reading the result.} Table~\ref{tab:compare} and
Fig.~\ref{fig:compare} make the case quantitatively. The GA improves on the
baseline by \SI{27.3}{\percent}, clearly the best of the four. Pure random
search barely moves the needle (\SI{-8.8}{\percent}): the search space is far
too large to sample blindly. Hill climbing does much better
(\SI{-19.7}{\percent}) by exploiting local structure, but its greedy chains get
stuck in local optima. Simulated annealing, on this budget and neighbourhood,
spends too many of its evaluations accepting worsening moves and never catches
up. The GA combines the strengths the others lack --- a population that explores
broadly, crossover that recombines good per-intersection timings, and elitism
that locks in progress --- which is exactly why it reaches the deepest
improvement. The comparison thus turns ``we chose a GA'' from an assertion into
a measured conclusion.

% ====================================================================
\section{Parallelisation}
% ====================================================================
Each generation requires one SUMO run per plan, and those runs are completely
independent. This is the parallelism the system exploits. With a master--worker
pattern, the master (MPI rank 0) holds the GA state and scatters the
generation's plans to $w$ worker ranks; each worker runs SUMO on its plans and
returns the delays; the master gathers them and proceeds.

The achievable speedup is bounded by Amdahl's law. If a fraction $f$ of the work
is inherently serial (here: GA bookkeeping, scatter/gather, and the per-batch
load imbalance), the speedup on $w$ workers cannot exceed
\begin{equation}
  S(w) \;=\; \frac{1}{f + \dfrac{1 - f}{w}}.
  \label{eq:amdahl}
\end{equation}
Because evaluation dominates runtime, $f$ is small and the system stays close to
linear for moderate worker counts, with efficiency tapering as $w$ grows and the
fixed serial cost becomes proportionally larger (Sec.~\ref{sec:strong}).

% ====================================================================
\section{Experimental setup}
% ====================================================================
All HPC measurements were taken on \textbf{Cyfronet Ares} (PLGrid grant
\texttt{plglscclass26-cpu}, partition \texttt{plgrid}), on a single compute node
with up to 48 MPI ranks. SUMO v1.20.0 was \emph{compiled from source} on the
cluster with the \texttt{libsumo} in-process bindings, because the pre-built
\texttt{eclipse-sumo} wheel segfaults on Ares compute nodes and ships only the
slower socket-based \texttt{traci}. OpenMPI 4.1.6 (matching the GCC 13.2.0 used
to build SUMO) launches the ranks, with \texttt{SLURM\_MPI\_TYPE=pmix} so that
\texttt{srun} starts them via PMIx.

\paragraph{How a study is run.} Each study is a self-contained SLURM batch
script in \texttt{scripts/slurm/}, submitted with \texttt{sbatch}. The script
loads the cluster modules, activates the \texttt{traffic} Conda environment,
exports the source-built \texttt{SUMO\_HOME}, and launches the parallel run:
\begin{verbatim}
module load miniconda3/24.5.0-0
module load openmpi/4.1.6-gcc-13.2.0   # libmpi + PMIx for srun-launched ranks
conda activate traffic
export SUMO_HOME=$HOME/sumo-1_20_0; export SLURM_MPI_TYPE=pmix
srun python -m mpi4py.futures -m src.ga.parallel  ...  # the parallel GA
\end{verbatim}
The \texttt{-m mpi4py.futures} launcher starts an \texttt{MPIPoolExecutor}:
\textbf{rank 0 is the master} that holds the GA state, and the remaining
$w = (\text{ranks}-1)$ ranks are \textbf{workers}. Every generation the master
scatters the population's plans across the workers, each worker runs SUMO
in-process on its plans and returns the delays, and the master gathers them ---
the master--worker pattern of the parallelisation section. Each job appends its
measurements to a CSV on the node; the plots in this report are rendered
\emph{locally} from those CSVs (Matplotlib is deliberately not installed on the
cluster, so a job's run is never blocked by plotting).

\paragraph{How the variables were swept.}
\begin{itemize}
\item \emph{Strong scaling} and \emph{throughput} (Tables~\ref{tab:strong},
\ref{tab:throughput}): the city is fixed at $8\times8$ and the run is repeated
for worker counts $w \in \{1,3,7,15,23,31,47\}$, so the same problem is solved
with a growing pool.
\item \emph{Problem size} (Table~\ref{tab:size}): the worker count is fixed at
$w=23$ and the city is swept from $3\times3$ (9 intersections, $n=18$) to
$32\times32$ (1024 intersections, $n=2048$).
\end{itemize}
The interactive dashboard demonstration in Sec.~\ref{sec:dashboard} instead uses
a smaller $3\times3$ grid so the optimisation completes in seconds on a laptop;
it is a usability demo, not an HPC measurement.

% ====================================================================
\section{Results}
% ====================================================================
\subsection{Optimisation quality}
On the $8\times8$ city the GA reduces total vehicle delay from
\SI{39473}{\veh\second} (fixed-time) to \SI{28689}{\veh\second}, an improvement
of \textbf{\SI{27.3}{\percent}}. The same optimised plan is produced regardless
of worker count (Table~\ref{tab:strong}), confirming that parallelisation
changes only the time-to-solution, not the solution itself.

\subsection{Strong scaling}
\label{sec:strong}
Here the problem is fixed ($8\times8$ city, population 128, 20 generations) and
the worker count is increased. Table~\ref{tab:strong} and
Fig.~\ref{fig:strong} report the wall time, the resulting speedup
$S(w)=t_1/t_w$ and the parallel efficiency $E(w)=S(w)/w$.

\begin{table}[H]
\centering
\caption{Strong scaling on the $8\times8$ city (fixed problem, growing workers).}
\label{tab:strong}
\begin{tabular}{S[table-format=2.0]S[table-format=4.2]S[table-format=2.2]S[table-format=1.3]}
\toprule
{Workers $w$} & {Wall time [s]} & {Speedup $S$} & {Efficiency $E$} \\
\midrule
1  & 1291.69 & 1.00  & 1.000 \\
3  & 437.04  & 2.96  & 0.985 \\
7  & 199.51  & 6.47  & 0.925 \\
15 & 100.61  & 12.84 & 0.856 \\
23 & 72.01   & 17.94 & 0.780 \\
31 & 60.09   & 21.50 & 0.694 \\
47 & 47.92   & 26.96 & 0.574 \\
\bottomrule
\end{tabular}
\end{table}

\begin{figure}[H]
\centering
\includegraphics[width=0.85\textwidth]{strong_scaling.png}
\caption{Strong scaling: speedup and parallel efficiency versus worker count.
Speedup stays near-linear up to $\sim$15 workers and reaches $26.96\times$ on
47 workers; efficiency declines as the fixed serial fraction of
Eq.~\eqref{eq:amdahl} grows in relative weight.}
\label{fig:strong}
\end{figure}

The behaviour is textbook Amdahl. Up to 15 workers the system is over
\SI{85}{\percent} efficient; beyond that, the per-generation serial overhead
(scatter/gather plus the wait for the slowest worker in each batch) becomes a
larger share of an ever-shorter runtime, so efficiency tapers to
\SI{57}{\percent} at 47 workers. The practically interesting point is the
time-to-solution: a run that takes \SI{21.5}{\minute} serially completes in
\SI{48}{\second} on 47 workers.

\subsection{Throughput}
Throughput measures raw evaluation rate: how many SUMO runs per second the pool
sustains as workers are added (fixed $8\times8$ city, 480 evaluations).

\begin{table}[H]
\centering
\caption{Throughput on the $8\times8$ city (\texttt{libsumo} backend).}
\label{tab:throughput}
\begin{tabular}{S[table-format=2.0]S[table-format=2.3]S[table-format=3.1]}
\toprule
{Workers $w$} & {Evaluations / s} & {ms / evaluation} \\
\midrule
1  & 1.447  & 691.1 \\
3  & 4.339  & 230.5 \\
7  & 9.981  & 100.2 \\
15 & 20.359 & 49.1 \\
23 & 28.940 & 34.6 \\
31 & 36.103 & 27.7 \\
47 & 45.776 & 21.8 \\
\bottomrule
\end{tabular}
\end{table}

\begin{figure}[H]
\centering
\includegraphics[width=0.85\textwidth]{throughput.png}
\caption{Throughput: sustained SUMO evaluations per second against worker count.
The rate rises from \SI{1.4}{evaluations\per\second} on a single worker to
\SI{45.8}{evaluations\per\second} on 47 workers --- a $31\times$ increase in raw
evaluation capacity.}
\label{fig:throughput}
\end{figure}

Throughput climbs from \SI{1.4}{evaluations\per\second} (one worker,
\SI{691}{\milli\second} per simulation) to \SI{45.8}{evaluations\per\second} on
47 workers. The in-process \texttt{libsumo} backend is what makes each
evaluation cheap enough for this to matter; a socket-based backend would spend a
large fixed cost per run and cap the achievable rate.

\subsection{Problem size}
Finally we hold the worker count fixed (23) and grow the city, measuring how the
cost per evaluation scales with the number of intersections.

\begin{table}[H]
\centering
\caption{Cost per evaluation versus city size (23 workers, \texttt{libsumo}).}
\label{tab:size}
\begin{tabular}{S[table-format=2.0]S[table-format=4.0]S[table-format=4.0]S[table-format=2.3]S[table-format=3.1]}
\toprule
{Grid} & {Intersections} & {Genome $n$} & {Evaluations / s} & {ms / eval} \\
\midrule
3  & 9    & 18   & 37.048 & 27.0  \\
5  & 25   & 50   & 35.969 & 27.8  \\
8  & 64   & 128  & 21.383 & 46.8  \\
12 & 144  & 288  & 12.815 & 78.0  \\
16 & 256  & 512  & 9.022  & 110.8 \\
24 & 576  & 1152 & 6.090  & 164.2 \\
32 & 1024 & 2048 & 4.730  & 211.4 \\
\bottomrule
\end{tabular}
\end{table}

\begin{figure}[H]
\centering
\includegraphics[width=0.85\textwidth]{problem_size.png}
\caption{Problem size: evaluation rate and per-evaluation cost as the city grows
from 9 to 1024 intersections. The cost per evaluation rises sub-quadratically
with city size, so larger cities remain tractable on the same worker pool.}
\label{fig:size}
\end{figure}

A $32\times32$ city has $113\times$ more intersections than a $3\times3$ grid,
yet the cost per evaluation grows only about $8\times$ (from \SI{27}{\milli
\second} to \SI{211}{\milli\second}). The simulator scales gracefully with the
network, so optimising realistically large cities is a matter of adding workers,
not of hitting an algorithmic wall.

% ====================================================================
\section{The interactive dashboard}
\label{sec:dashboard}
% ====================================================================
The dashboard is the project's centrepiece: it lets a user build a network,
optimise its signals, and \emph{see} the effect, without touching the command
line. The screenshots below are from a $3\times3$ demonstration grid (baseline
\SI{15638}{\veh\second}, optimised \SI{12124}{\veh\second},
\SI{-22.5}{\percent}).

\begin{figure}[H]
\centering
\includegraphics[width=0.92\textwidth]{shot_network.png}
\caption{\textbf{Network builder.} The left panel parametrises the city ---
topology, junctions per side, lanes, edge length, demand intensity and the
simulation horizon. Clicking \emph{Build \& show network} generates the SUMO
network and renders it. Here a $3\times3$ grid with single-lane roads and a
vehicle every \SI{1.6}{\second} is created.}
\label{fig:shot-network}
\end{figure}

\begin{figure}[H]
\centering
\includegraphics[width=0.92\textwidth]{shot_convergence.png}
\caption{\textbf{Optimisation result.} The headline cards report baseline,
optimised and improvement (\SI{-22.5}{\percent}). \emph{Left:} the GA
convergence curve --- total delay drops generation by generation and plateaus,
indicating the search has converged. \emph{Right:} the same before/after delay
as bars.}
\label{fig:shot-convergence}
\end{figure}

\begin{figure}[H]
\centering
\includegraphics[width=0.78\textwidth]{shot_animation.png}
\caption{\textbf{Traffic animation.} The optimised plan is replayed: each dot is
a vehicle, coloured by speed (green = moving, red = halting). A time slider
scrubs through the simulated horizon, making queue formation and clearance
visible directly.}
\label{fig:shot-animation}
\end{figure}

\begin{figure}[H]
\centering
\includegraphics[width=0.92\textwidth]{shot_heatmap.png}
\caption{\textbf{Congestion heatmap.} Baseline (left) versus optimised (right),
with each intersection coloured by its average number of halting vehicles
(green = clear, red = queued). The optimised network shows visibly fewer hot
spots --- the same \SI{-22.5}{\percent} improvement, now localised in space.}
\label{fig:shot-heatmap}
\end{figure}

The dashboard also supports a Cyfronet round-trip: a network can be exported as
a bundle, trained on Ares with the parallel GA, and the resulting plan imported
back to redraw the before/after charts --- closing the loop between the
single-machine front-end and the HPC back-end.

% ====================================================================
\section{Discussion}
% ====================================================================
Three threads come together in this project. First, framing signal timing as
black-box minimisation of simulated delay (Eq.~\eqref{eq:delay}) makes the
problem precise and lets a generic optimiser attack any network the dashboard
can build. Second, the algorithm comparison (Sec.~\ref{sec:compare}) shows the
GA is not an arbitrary choice: on an equal budget it beats random search, hill
climbing and simulated annealing, reaching \SI{-27.3}{\percent} where the others
reach \SI{-19.7}{\percent} or less. Third, the independence of fitness
evaluations turns the optimiser into an HPC workload that scales to a
$26.96\times$ speedup and \SI{45.8}{evaluations\per\second} on Ares, while the
problem-size study confirms the simulator itself scales gracefully to
thousand-intersection cities.

The efficiency decline at high worker counts (Sec.~\ref{sec:strong}) is the
expected consequence of Amdahl's law: with evaluation so cheap under
\texttt{libsumo}, the fixed per-generation serial cost eventually dominates.
This is a feature of the regime, not a defect --- the practically relevant
metric, time-to-solution, still falls from over twenty minutes to under a minute.

% ====================================================================
\section{Conclusions}
% ====================================================================
We built a complete, reproducible pipeline --- dashboard, genetic algorithm and
SUMO simulator --- that optimises urban traffic-signal timings and scales the
optimisation on a supercomputer. The optimised plans cut total vehicle delay by
\SI{27.3}{\percent} on an $8\times8$ city and \SI{22.5}{\percent} on the
dashboard demonstration, all by retiming existing green phases without any new
infrastructure. An equal-budget comparison against three alternative search
strategies confirms that the genetic algorithm is the right tool for this noisy,
gradient-free objective. By dispatching every fitness evaluation through an
abstract parallel map, the very same optimiser runs unchanged on a laptop and on
Cyfronet Ares, where it achieves a $26.96\times$ speedup on 47 workers and
sustains up to \SI{45.8}{SUMO evaluations\per\second}. Together these results
show that high-performance computing makes data-driven signal optimisation
practical at city scale: the bottleneck is no longer the search, but simply the
number of workers one is willing to spend.

\end{document}
